\documentclass[12pt,oneside]{book}

% ============================================================
% METADATOS
% ============================================================
\newcommand{\universidad}{Universidad de Buenos Aires (UBA)}
\newcommand{\departamento}{C\'atedra de Pr\'actica Profesional Supervisada\\[0.15cm] \'Area Derecho de Familia y Procesal}
\newcommand{\materia}{Pr\'actica Profesional en Derecho}
\newcommand{\titulo}{Cl\'inica Jur\'idica --- Consultorio Gratuito}
\newcommand{\autor}{Isaac Edgar Camacho Ocampo}
\newcommand{\anio}{2026}

% ============================================================
% PAQUETES
% ============================================================
\usepackage[utf8]{inputenc}
\usepackage[T1]{fontenc}
\usepackage[spanish,es-nodecimaldot]{babel}

\usepackage[
    top=2.5cm,
    bottom=2.5cm,
    left=3cm,
    right=2.5cm,
    headheight=15pt
]{geometry}

\usepackage{amsmath}
\usepackage{amssymb}
\usepackage{mathtools}

\usepackage{graphicx}
\usepackage{float}
\usepackage{caption}
\usepackage{subcaption}

\usepackage{booktabs}
\usepackage{array}
\usepackage{longtable}
\usepackage{tabularx}

\usepackage{enumitem}

\usepackage{xcolor}
\usepackage[most]{tcolorbox}

\usepackage{fancyhdr}
\usepackage{ragged2e}

\usepackage[
    colorlinks=true,
    linkcolor=blue!50!black,
    citecolor=green!40!black,
    urlcolor=blue!60!black,
    pdftitle={\titulo},
    pdfauthor={\autor},
    pdfsubject={Apuntes universitarios},
    bookmarks=true
]{hyperref}

% ============================================================
% CONFIGURACIÓN
% ============================================================
\graphicspath{
    {./img/}
    {./graficos/}
    {./figuras/}
}

\setcounter{tocdepth}{2}

\setlength{\parindent}{0pt}
\setlength{\parskip}{0.5em}

\setlist[itemize]{
    topsep=0.3em,
    itemsep=0.2em
}

\setlist[enumerate]{
    topsep=0.3em,
    itemsep=0.2em
}

\clubpenalty=10000
\widowpenalty=10000

% ============================================================
% COMANDOS
% ============================================================
\newcommand{\abs}[1]{\left|#1\right|}
\newcommand{\norm}[1]{\left\lVert#1\right\rVert}
\newcommand{\vect}[1]{\mathbf{#1}}
\newcommand{\deriv}[2]{\frac{d#1}{d#2}}
\newcommand{\pderiv}[2]{\frac{\partial#1}{\partial#2}}

% ============================================================
% ENTORNOS / CAJAS
% ============================================================
\newtcolorbox{definition}[1]{
    enhanced,
    breakable,
    title={Definici\'on: #1},
    fonttitle=\bfseries,
    colback=blue!3,
    colframe=blue!50!black,
    boxrule=0.8pt,
    arc=2mm
}

\newtcolorbox{important}{
    enhanced,
    breakable,
    title={Importante},
    fonttitle=\bfseries,
    colback=yellow!8,
    colframe=orange!70!black,
    boxrule=0.8pt,
    arc=2mm
}

\newtcolorbox{example}[1]{
    enhanced,
    breakable,
    title={Ejemplo: #1},
    fonttitle=\bfseries,
    colback=green!4,
    colframe=green!40!black,
    boxrule=0.8pt,
    arc=2mm
}

\newtcolorbox{formula}[1]{
    enhanced,
    breakable,
    title={F\'ormula / Regla: #1},
    fonttitle=\bfseries,
    colback=purple!4,
    colframe=purple!50!black,
    boxrule=0.8pt,
    arc=2mm
}

\newtcolorbox{exercise}[1]{
    enhanced,
    breakable,
    title={Ejercicio: #1},
    fonttitle=\bfseries,
    colback=gray!5,
    colframe=gray!60!black,
    boxrule=0.8pt,
    arc=2mm
}

\newtcolorbox{note}{
    enhanced,
    breakable,
    title={Nota},
    fonttitle=\bfseries,
    colback=white,
    colframe=gray!50,
    boxrule=0.6pt,
    arc=2mm
}

% ============================================================
% CONFIGURACIÓN FINAL
% ============================================================
\pagestyle{fancy}
\fancyhf{}
\fancyhead[L]{\nouppercase{\leftmark}}
\fancyhead[R]{\materia}
\fancyfoot[C]{\thepage}
\renewcommand{\headrulewidth}{0.4pt}
\renewcommand{\footrulewidth}{0pt}

\fancypagestyle{plain}{
    \fancyhf{}
    \fancyfoot[C]{\thepage}
    \renewcommand{\headrulewidth}{0pt}
}

\renewcommand{\arraystretch}{1.25}

% ============================================================
% DOCUMENTO
% ============================================================
\begin{document}

% ------------------------------------------------------------
% PORTADA
% ------------------------------------------------------------
\begin{titlepage}

\thispagestyle{empty}

\begin{center}

\vspace*{1cm}

{\large \textsc{\universidad}}

\vspace{0.4cm}

{\large \departamento}

\vspace{3cm}

{\Huge \textbf{\materia}}

\vspace{1cm}

{\LARGE \titulo}

\vspace{0.8cm}

\textit{Comprender lo que se estudia es el primer paso para transformarlo en conocimiento propio.}

\vspace{1.2cm}

\rule{0.70\textwidth}{0.5pt}

\vspace{0.5cm}

{\large Apuntes de estudio}

\vspace{0.5cm}

\rule{0.70\textwidth}{0.5pt}

\vfill

{\large \autor}

\vspace{0.3cm}

{\large \anio}

\end{center}

\vfill

\begin{flushleft}
\textit{Este es un modesto aporte a los estudiantes de ``Pr\'actica Profesional en Derecho'' no pretende sustituir el material oficial de la materia.}
\end{flushleft}

\end{titlepage}

% ------------------------------------------------------------
% ÍNDICE
% ------------------------------------------------------------
\tableofcontents

% ------------------------------------------------------------
% PRESENTACIÓN
% ------------------------------------------------------------
\chapter*{Presentaci\'on}
\addcontentsline{toc}{chapter}{Presentaci\'on}
\markboth{Presentaci\'on}{Presentaci\'on}

Estos apuntes documentan una sesi\'on real de trabajo en un consultorio jur\'idico gratuito de pr\'actica profesional. A diferencia de una clase te\'orica tradicional, se registra el funcionamiento efectivo de una cl\'inica jur\'idica: la atenci\'on a una consultante por un caso de familia (cuota alimentaria, r\'egimen de comunicaci\'on, embargo de haberes y mediaci\'on), la gesti\'on administrativa interna del equipo (expedientes, notificaciones, cuentas bancarias) y, hacia el final, la orientaci\'on brindada a los alumnos que ingresan por primera vez sobre la obligatoriedad de asistir a audiencias y el sentido profesional de la pr\'actica.

\textbf{Utilidad para la vida profesional.} El material muestra el derecho tal como se ejerce y no solo como se estudia: un cliente real que necesita resultados concretos, expedientes que se atrasan, domicilios que no se ubican, bancos que descuentan sumas por error, compa\~neros que coordinan tareas por chat y un profesional que decide cu\'ando hablar y cu\'ando escuchar. Ser abogado es, en gran medida, gestionar informaci\'on, coordinar personas y sostener la paciencia frente a la lentitud del sistema.

\textbf{Utilidad para la vida personal.} El caso de familia expone c\'omo el desacuerdo entre dos adultos repercute en la vida cotidiana de una ni\~na, y c\'omo la falta de acuerdos claros (de dinero, de tiempos, de comunicaci\'on) genera desgaste emocional sostenido. Comprender ese mecanismo ayuda a poner por escrito lo que se promete, a no delegar en el silencio lo que deber\'ia hablarse y a reconocer que la falta de previsi\'on de hoy es el conflicto de ma\~nana.

\begin{note}
El audio original corresponde al registro ambiente de un consultorio jur\'idico gratuito y no a una clase magistral. Contiene tramos con superposici\'on de voces, frases incompletas y segmentos que la transcripci\'on autom\'atica no logr\'o captar con claridad; esos tramos se indican como \emph{[tramo inaudible]}. No se incorporan art\'iculos de c\'odigo, ley o normativa porque ninguno de los intervinientes los cit\'o de forma textual y verificable durante la grabaci\'on. Por indicaci\'on de la consigna original, no se menciona el nombre de la instituci\'on ni de las personas docentes o consultantes; se las identifica por su rol. La fecha de la grabaci\'on no figura en el registro original.
\end{note}

% TODO: contenido incompleto o ambiguo en las notas originales (fecha de la sesion no consignada)

\begin{center}
\small
\begin{tabularx}{\textwidth}{>{\raggedright\arraybackslash}p{0.28\textwidth}>{\raggedright\arraybackslash}X}
\toprule
\textbf{Bloque} & \textbf{Contenido} \\
\midrule
1. Atenci\'on al consultante & Consulta inicial, sost\'en econ\'omico de la hija, r\'egimen de comunicaci\'on y cuidado. \\
2. Gesti\'on del proceso & Mediaci\'on previa, notificaci\'on al demandado no ubicable, embargo de haberes. \\
3. Funcionamiento interno & Coordinaci\'on de equipo y expedientes; tr\'amites bancarios y administrativos. \\
4. Formaci\'on del practicante & Audiencias, actitud profesional y rol pr\'actico del abogado. \\
5. Cierre & Compromiso social de la cl\'inica. \\
\bottomrule
\end{tabularx}
\end{center}

% ============================================================
% CAPÍTULO 1
% ============================================================
\chapter[Atenci\'on al consultante]{Atenci\'on al consultante: caso de familia}

La sesi\'on se abre con la atenci\'on de una consultante que expone, junto a un pasante del consultorio, la situaci\'on de incumplimiento de la cuota alimentaria por parte del padre de su hija. A lo largo de este bloque se reconstruye la din\'amica econ\'omica y afectiva del conflicto familiar.

\section{Consulta inicial y contexto del caso}

El intercambio comienza con una pregunta de contexto sobre la finalidad de un documento solicitado por la consultante. El pasante aclara que el tr\'amite puede realizarse con la documentaci\'on disponible y que lo relativo al caso consta en el oficio; ante cualquier duda, la consultante debe ponerse en comunicaci\'on con la persona correspondiente.
% TODO: contenido incompleto o ambiguo en las notas originales (el pasante no completa con quien debe comunicarse la consultante)

La consultante explica que existe un n\'umero de expediente o comunicaci\'on vinculado al caso, y que la contraparte (el padre de la ni\~na) habr\'ia anticipado verbalmente que \textit{``el otro mes le van a embargar''}.
% TODO: contenido incompleto o ambiguo en las notas originales (no se aclara a quien se refiere el embargo anticipado)

Relata que el padre le enviaba numerosos audios y que, en uno de ellos, ella le respondi\'o que, si realmente quer\'ia conocer su domicilio, lo conseguir\'ia en dos minutos. Indica adem\'as que el padre \textbf{dej\'o de pagar el colegio}, hecho que la consultante confirma como efectivamente ocurrido.

La consultante expone su lectura de la conducta de la contraparte: entiende que \'esta buscar\'ia que se le generara una \textbf{deuda} para que, al mes siguiente, cuando ella deba pagar con el dinero que \'el le transferir\'a (porque no dejar\'a de pagar el colegio), el monto adeudado se encuentre incrementado. Se menciona una suma total de cuarenta y ocho (sin que se llegue a precisar la cifra completa).
% TODO: contenido incompleto o ambiguo en las notas originales (el monto total queda trunco; tramo inaudible)

\section{Sost\'en econ\'omico y gastos de la hija}

La consultante detalla que se encuentra costeando por s\'i misma el \textbf{viaje de egresados} y la campera correspondiente, y solicita que se le remitan los comprobantes de las cuotas para acreditar los pagos realizados, incluida una transferencia ya efectuada.

Relata un episodio de enojo de la contraparte al advertir que ella solventaba el pago de la remera de egresados sin haberlo consultado previamente. Su fundamento es el siguiente: si el padre paga el colegio con dos o tres meses de demora, depender de \'el para un viaje, cuyas cuotas son mensuales, resultar\'ia inviable; adem\'as, desea evitar tener que explicar a otros padres del colegio que el padre paga cuando quiere. Sostiene tambi\'en que tiene derecho a saber que el padre est\'a pagando el viaje, y se pregunta c\'omo \'el podr\'ia saber si ella iba a querer colaborar.

La \textbf{irregularidad de los aportes} se ejemplifica con dos pedidos concretos que no obtuvieron respuesta: diez mil pesos para una prenda puntual y dos mil pesos para una merienda. La consultante se\~nala que el problema es constante y que, ni siquiera ante pedidos menores, recibe respuesta.

Cierra el punto anunciando su decisi\'on de \textbf{no continuar intercediendo por la contraparte} ante eventuales medidas de embargo: que se le retire lo que corresponda y que, si quiere o no pagar el colegio, es asunto suyo.

Se menciona tambi\'en que el colegio hab\'ia otorgado a la ni\~na un \textbf{beneficio educativo} (ayuda parcial) de aproximadamente veinte mil pesos en un ejercicio anterior, en atenci\'on a la irregularidad de los pagos; la instituci\'on habr\'ia indicado que dicho beneficio reci\'en podr\'ia renovarse al a\~no siguiente.

\section{R\'egimen de comunicaci\'on y cuidado personal}

La consultante describe la inestabilidad econ\'omica que le genera no poder prever los aportes de la contraparte, lo que le impide sostener gastos regulares para la ni\~na: si en un mes compra los pantalones, el mes siguiente no puede comprar las zapatillas. Se maneja sola, porque el padre ``est\'a o no est\'a''.

Respecto del r\'egimen de contacto, se menciona un esquema de \textbf{fin de semana por medio}. La consultante expresa preocupaci\'on por la falta de supervisi\'on adecuada durante esos per\'iodos y remarca que el verdadero objetivo del r\'egimen de comunicaci\'on deber\'ia ser \textbf{sostener el v\'inculo entre padre e hija} y no \'unicamente el aspecto econ\'omico. Indica que la ni\~na no manifiesta ganas de ir a lo de su padre: sabe que debe ir porque la madre trabaja y, si la abuela materna no puede, no hay quien la cuide. Seg\'un expresa, la raz\'on no es econ\'omica.

Sobre su necesidad de generar ingresos, explica que trabajando medio tiempo percibe quinientos ochenta mil pesos, suma que considera insuficiente. La contraparte le habr\'ia respondido que trabaje m\'as o busque otro empleo; la consultante replica que no tiene qui\'en cuide a la ni\~na y que no puede costear una ni\~nera.

Finalmente, describe una modalidad de trato que caracteriza como \textbf{intimidatoria}, vinculada a la obra social de la ni\~na y a amenazas veladas de no buscarla ni cuidarla, lo que, seg\'un la consultante, le genera todas las consecuencias a ella.

\begin{important}
Seg\'un la consultante, el objetivo central del r\'egimen de comunicaci\'on no es econ\'omico: debe sostener el v\'inculo entre padre e hija y quedar \textbf{pactado}, de modo que no dependa de la voluntad variable de la contraparte.
\end{important}

\section*{Cuadro Cornell: Atenci\'on al consultante}
\addcontentsline{toc}{section}{Cuadro Cornell: Atenci\'on al consultante}

\begin{table}[H]
\centering
\small
\begin{tabularx}{\textwidth}{
    >{\raggedright\arraybackslash}p{0.28\textwidth}
    >{\raggedright\arraybackslash}X
}
\toprule
\textbf{Preguntas / palabras clave} &
\textbf{Notas / respuestas} \\
\midrule

\textbf{¿Qu\'e conductas de la contraparte se relatan en el caso de familia?} &
Dej\'o de pagar el colegio, no respondi\'o pedidos de dinero menores (diez mil pesos para una prenda, dos mil para una merienda), enoj\'o ante el pago por la consultante de la remera de egresados y amenaz\'o con no buscar ni cuidar a la ni\~na. \\

\textbf{¿Qu\'e lectura hace la consultante de la deuda por el colegio?} &
Entiende que la contraparte busca que se le genere una deuda para que, el mes siguiente, el monto a pagar con la transferencia que \'el le haga se vea incrementado. \\

\textbf{¿Qu\'e gastos costea la consultante por s\'i misma?} &
El viaje de egresados y la campera. Solicita comprobantes de cuota para acreditar los pagos, incluida una transferencia ya realizada. \\

\textbf{¿Por qu\'e la consultante decidi\'o no intervenir m\'as ante un eventual embargo?} &
Porque considera que si la contraparte quiere o no pagar el colegio es su problema; anuncia que no la defender\'a m\'as. \\

\textbf{¿Qu\'e beneficio otorg\'o el colegio y qu\'e se indic\'o sobre su renovaci\'on?} &
Un beneficio educativo de aproximadamente veinte mil pesos, otorgado en atenci\'on a la irregularidad de los pagos; la instituci\'on habr\'ia indicado que solo podr\'ia renovarse el a\~no siguiente. \\

\textbf{¿Qu\'e r\'egimen de contacto se menciona y qu\'e objetivo debe tener?} &
Fin de semana por medio. Seg\'un la consultante, su objetivo debe ser sostener el v\'inculo entre padre e hija, y no solo el aspecto econ\'omico. \\

\textbf{¿Qu\'e dificultades de organizaci\'on plantea la consultante?} &
No puede prever los aportes de la contraparte ni sostener gastos regulares; trabaja medio tiempo (quinientos ochenta mil pesos), no tiene qui\'en cuide a la ni\~na ni recursos para una ni\~nera, y la contraparte le indic\'o que trabaje m\'as. \\

\midrule

\multicolumn{2}{>{\raggedright\arraybackslash}X}{
\textbf{Resumen:}
El caso expone un incumplimiento sostenido de aportes alimentarios (colegio, gastos de egreso, pedidos menores) que genera inestabilidad econ\'omica en la consultante. A ello se suma la preocupaci\'on por el cuidado de la ni\~na durante los fines de semana y por el trato intimidatorio de la contraparte. La consultante busca un r\'egimen de comunicaci\'on pactado que sostenga el v\'inculo padre--hija.
} \\

\bottomrule
\end{tabularx}
\end{table}

% ============================================================
% CAPÍTULO 2
% ============================================================
\chapter[Gesti\'on del proceso]{Gesti\'on del proceso: mediaci\'on, notificaci\'on y embargo}

\section{Mediaci\'on previa}

El pasante informa a la consultante que se hab\'ia concurrido previamente para gestionar una \textbf{fecha de mediaci\'on} referida al r\'egimen de comunicaci\'on. El tr\'amite se realiza de forma \textbf{presencial} y la fecha se otorga en el momento. Una vez generado el tr\'amite, el equipo debe llevar la documentaci\'on ya firmada por la consultante; all\'i se otorga la fecha y la contraparte es notificada, incluso por \textbf{carta documento}, para la mediaci\'on.

La consultante explica qu\'e espera obtener del proceso: un compromiso concreto y verificable sobre los d\'ias de cuidado, que no quede sujeto a la voluntad discrecional de la contraparte (por ejemplo, un d\'ia determinado o semana por medio, pero \textbf{que quede pactado}). Se\~nala que actualmente la contraparte puede indicar un d\'ia que no puede o no quiere cuidar a la ni\~na, y ella debe ingeni\'arselas.

Agrega que su disponibilidad laboral es relativamente flexible entre semana, pero no los fines de semana, por la carga de trabajo propia de su empleo (mencionado en referencia a un centro comercial), y que para esos d\'ias recurre a la colaboraci\'on de su madre.

\section{Notificaci\'on al demandado no ubicable}

En un segundo caso, de otra consultante, se aborda la dificultad para notificar a la contraparte porque el domicilio denunciado presenta inconsistencias: el departamento indicado (n\'umero nueve) no existir\'ia en el edificio, cuya numeraci\'on llegar\'ia solamente hasta el cinco. La c\'edula volvi\'o negativa por este motivo.

Se plantea como estrategia procesal invocar la \textbf{demora} que implicar\'ia un nuevo libramiento de c\'edula a la provincia (estimada en quince d\'ias, como la anterior) sumada a la demora adicional de la jurisdicci\'on provincial, como fundamento para solicitar una v\'ia de notificaci\'on alternativa.

Se eval\'ua tambi\'en notificar por \textbf{WhatsApp} ``de manera excepcional'', y se analiza como alternativa pr\'actica ubicar a una t\'ia de la contraparte, a cuyo domicilio se habr\'ian hecho llegar previamente env\'ios por plataformas de pago y mensajer\'ia, como v\'ia indirecta m\'as confiable para lograr el contacto: all\'i s\'i se recibir\'ia la notificaci\'on.

Se hace referencia, adem\'as, a que el n\'umero de tel\'efono de la contraparte se encuentra corroborado por figurar denunciado en el expediente y en la planilla de la Oficina de Violencia Dom\'estica (OVD).

\begin{important}
Cuando la c\'edula resulta negativa por inexistencia del domicilio, las v\'ias discutidas por el equipo fueron: argumentar la demora de un nuevo libramiento a la provincia, notificar excepcionalmente por WhatsApp (n\'umero corroborado en el expediente y en la planilla de la OVD) y ubicar el domicilio de una familiar de la contraparte.
\end{important}

\section{Medida cautelar: embargo de haberes}

Este segmento documenta la coordinaci\'on interna para librar un \textbf{oficio de embargo} sobre una cuenta bancaria de la contraparte. Se discute qu\'e banco corresponde notificar, dado que exist\'ian dos entidades o dos personas con nombre coincidente, lo que gener\'o confusi\'on sobre el domicilio y la jurisdicci\'on correctos (capital o provincia): uno de los homónimos era de capital y el otro de provincia. El domicilio proporcionado era de provincia, y el otro aparec\'ia tambi\'en de provincia con un n\'umero de expediente distinto, por lo que se hab\'ia intuido que era este \'ultimo.
% TODO: contenido incompleto o ambiguo en las notas originales (identificacion de la cuenta a individualizar: tramo inaudible)

Se hace referencia a la f\'ormula procesal empleada para solicitar el embargo mediante transferencia y oficio, en los t\'erminos pronunciados durante la grabaci\'on:

\begin{formula}{Solicitud de embargo por transferencia y oficio}
\textit{``Se manifiesta, se manifiesta, se manifiesta. Lo solicito por transferencia y se libre el oficio.''}
\end{formula}

En un tercer caso se relata un descuento bancario indebido sobre una cuenta sueldo por una deuda de mantenimiento de tarjeta con una antig\"uedad superior a cinco a\~nos. Se concluye que, al no existir movimientos que acreditar (la cuenta no se usaba), el equipo no puede intervenir m\'as all\'a de dejar constancia de un detalle de fechas.

\section*{Cuadro Cornell: Gesti\'on del proceso}
\addcontentsline{toc}{section}{Cuadro Cornell: Gesti\'on del proceso}

\begin{table}[H]
\centering
\small
\begin{tabularx}{\textwidth}{
    >{\raggedright\arraybackslash}p{0.28\textwidth}
    >{\raggedright\arraybackslash}X
}
\toprule
\textbf{Preguntas / palabras clave} &
\textbf{Notas / respuestas} \\
\midrule

\textbf{¿C\'omo se gestiona la mediaci\'on previa en el caso relatado?} &
De forma presencial, con la documentaci\'on firmada por la consultante; la fecha se otorga en el momento y la contraparte es notificada, incluso por carta documento. \\

\textbf{¿Qu\'e espera la consultante de la mediaci\'on?} &
Un compromiso concreto y verificable sobre los d\'ias de cuidado (d\'ia determinado o semana por medio) que quede pactado y no dependa de la voluntad discrecional de la contraparte. \\

\textbf{¿Qu\'e problema present\'o la notificaci\'on en el segundo caso?} &
La c\'edula volvi\'o negativa porque el departamento denunciado (n\'umero nueve) no existe: la numeraci\'on del edificio va del uno al cinco. \\

\textbf{¿Qu\'e v\'ias alternativas se evaluaron para notificar?} &
Invocar la demora de un nuevo libramiento a la provincia (quince d\'ias m\'as la demora provincial), notificar excepcionalmente por WhatsApp y ubicar a una t\'ia de la contraparte. \\

\textbf{¿Qu\'e respalda el n\'umero de tel\'efono de la contraparte?} &
Figura denunciado en el expediente y en la planilla de la Oficina de Violencia Dom\'estica (OVD). \\

\textbf{¿Qu\'e se debe verificar al librar un oficio de embargo?} &
La entidad bancaria y la jurisdicci\'on (capital o provincia), ya que en el caso hab\'ia dos homónimos que generaron confusi\'on. \\

\textbf{¿Qu\'e se concluy\'o ante el descuento bancario sobre la cuenta sueldo?} &
Que, al no existir movimientos que acreditar, el equipo solo puede dejar constancia de un detalle de fechas. \\

\midrule

\multicolumn{2}{>{\raggedright\arraybackslash}X}{
\textbf{Resumen:}
La gesti\'on del proceso combina la mediaci\'on previa (para pactar el r\'egimen de comunicaci\'on), la b\'usqueda de v\'ias alternativas cuando no se logra notificar a la contraparte y la tramitaci\'on del embargo mediante oficio y transferencia, verificando con cuidado entidad bancaria y jurisdicci\'on.
} \\

\bottomrule
\end{tabularx}
\end{table}

% ============================================================
% CAPÍTULO 3
% ============================================================
\chapter[Funcionamiento interno]{Funcionamiento interno del consultorio}

\section{Coordinaci\'on de equipo y expedientes}

Se documentan intercambios operativos entre los integrantes del equipo sobre el estado de distintos expedientes: notificaciones ya firmes, escritos pendientes de firma y la gesti\'on de un \textbf{mail centralizado} para la recepci\'on de datos personales de los nuevos integrantes.

Un integrante informa que una notificaci\'on ya fue practicada y que el asunto se encuentra firme; la fecha mencionada no se identifica con precisi\'on.
% TODO: contenido incompleto o ambiguo en las notas originales (fecha de la notificacion: tramo inaudible)

Otro integrante indica que un escrito ya est\'a redactado y que puede enviarse para su firma.

\section{Tr\'amites bancarios y administrativos}

Se relata, en primera persona, el episodio de una integrante del equipo sobre un \textbf{descuento bancario no reconocido}, vinculado a la mantenci\'on de una tarjeta asociada a una cuenta abierta con anterioridad. Aunque es ajeno al caso de la consultante principal, se conserva porque el docente lo utiliza como ejemplo de c\'omo verificar y documentar reclamos bancarios.

La integrante no hab\'ia pagado la mantenci\'on de la tarjeta porque no la usaba; la pas\'o a cuenta sueldo creyendo que no se la descontar\'ian, pero le descontaron doscientos mil pesos. El banco le indic\'o que, igualmente, exist\'ia una deuda y que le quedaban ochenta y nueve mil pesos por pagar, tras lo cual no se le descontar\'ia m\'as.

Se cierra el punto con la constataci\'on de que, al no existir controversia real sobre el uso de la cuenta, la v\'ia de reclamo institucional queda limitada a dejar asentado un detalle cronol\'ogico de fechas para un eventual reclamo posterior.
% TODO: contenido incompleto o ambiguo en las notas originales (el apunte presenta este episodio tanto como relato de una integrante del equipo como "tercer caso" del consultorio)

\section*{Cuadro Cornell: Funcionamiento interno}
\addcontentsline{toc}{section}{Cuadro Cornell: Funcionamiento interno}

\begin{table}[H]
\centering
\small
\begin{tabularx}{\textwidth}{
    >{\raggedright\arraybackslash}p{0.28\textwidth}
    >{\raggedright\arraybackslash}X
}
\toprule
\textbf{Preguntas / palabras clave} &
\textbf{Notas / respuestas} \\
\midrule

\textbf{¿Qu\'e aspectos se coordinan entre los integrantes del equipo?} &
El estado de los expedientes (notificaciones firmes, escritos pendientes de firma) y la gesti\'on de un mail centralizado para recibir los datos personales de los nuevos integrantes. \\

\textbf{¿Qu\'e episodio bancario se relata y por qu\'e se conserva?} &
Un descuento no reconocido por mantenci\'on de una tarjeta, de doscientos mil pesos, con una deuda remanente de ochenta y nueve mil pesos. Se conserva porque el docente lo usa como ejemplo de c\'omo verificar y documentar reclamos bancarios. \\

\textbf{¿Qu\'e puede hacer el equipo ante un reclamo bancario sin controversia sobre el uso de la cuenta?} &
Solo dejar asentado un detalle cronol\'ogico de fechas para un eventual reclamo posterior. \\

\midrule

\multicolumn{2}{>{\raggedright\arraybackslash}X}{
\textbf{Resumen:}
El funcionamiento del consultorio depende de la coordinaci\'on constante entre integrantes sobre expedientes, escritos y datos de los nuevos miembros, y del registro ordenado de fechas y comprobantes en los reclamos administrativos y bancarios.
} \\

\bottomrule
\end{tabularx}
\end{table}

% ============================================================
% CAPÍTULO 4
% ============================================================
\chapter[Formaci\'on del practicante]{Formaci\'on del practicante: audiencias y rol profesional}

Sobre el final de la sesi\'on, quien ejerce el rol docente se dirige a los alumnos que se incorporan por primera vez al consultorio, en el tramo m\'as claramente pedag\'ogico de toda la grabaci\'on.

\section{Obligatoriedad de asistir a audiencias}

Adem\'as de las tareas de escritorio, es \textbf{obligatorio} para los alumnos presenciar audiencias reales, tanto presenciales como virtuales. La asistencia constituye un elemento m\'as para aprobar la materia y debe concurrir todo el grupo. El docente fundamenta la exigencia en dos motivos.

\begin{enumerate}
    \item \textbf{Motivo pedag\'ogico:} permite observar c\'omo se realiza una audiencia verdadera, con personas reales.
    \item \textbf{Motivo pr\'actico:} los alumnos disponen de informaci\'on directa del caso que el equipo docente no tiene (si la parte pag\'o, si va a asistir). Los alumnos no est\'an ``de adorno'': son necesarios como parte del equipo.
\end{enumerate}

\begin{important}
Si ning\'un integrante del grupo puede asistir a una audiencia, debe informarse previamente al docente. Ante situaciones laborales, personales o familiares que impidan la concurrencia de todo el grupo, se eval\'ua cada caso, pero al menos una persona del grupo debe asistir: la ausencia total y no informada implica una calificaci\'on negativa para todo el grupo. Si los alumnos necesitan permiso en sus trabajos para asistir, el consultorio lo otorga.
\end{important}

\section{Reflexi\'on sobre una audiencia reciente}

El docente retoma una audiencia del d\'ia anterior para reflexionar sobre la actitud profesional observada en la contraparte, contrastando una postura colaborativa con la actitud habitual, m\'as apurada y desatenta, observada en otros procesos: los profesionales suelen tener muchas audiencias, salen al pasillo a conversar con los abogados sobre cu\'anto puede pagarse y cierran la audiencia casi sin escuchar a la persona. Seg\'un el docente, \textbf{la gente cree que es un tr\'amite}, y las audiencias lo corroboran.

Se\~nala, adem\'as, que en ocasiones el propio equipo debe moderar la extensi\'on de los relatos de sus consultantes durante la audiencia para no generar una impresi\'on desfavorable ante el secretario ni ante la contraparte. Recomienda a los alumnos estar cerca de sus clientes para poder, discretamente, interrumpir o distraer cuando sea necesario.

\section{Actitud profesional: la pr\'actica por sobre la teor\'ia}

El docente propone, como incentivo, llevar un registro informal (``tabla de posiciones'') de asistencia a tribunales, acreditada mediante el env\'io de una fotograf\'ia tomada en el lugar y subida al grupo. Habr\'a un premio y aclara que no es una broma.

Cierra este tramo con una definici\'on sobre el perfil profesional que busca formarse: \textbf{un abogado pr\'actico tiene m\'as posibilidades de resolver un caso que uno ``voltero''}, priorizando la resoluci\'on pr\'actica de los conflictos por sobre el apego estrictamente te\'orico. El docente indica que esta idea proviene de la pr\'actica.
% TODO: contenido incompleto o ambiguo en las notas originales (el termino "voltero" se conserva tal como figura en el apunte)

\section*{Cuadro Cornell: Formaci\'on del practicante}
\addcontentsline{toc}{section}{Cuadro Cornell: Formaci\'on del practicante}

\begin{table}[H]
\centering
\small
\begin{tabularx}{\textwidth}{
    >{\raggedright\arraybackslash}p{0.28\textwidth}
    >{\raggedright\arraybackslash}X
}
\toprule
\textbf{Preguntas / palabras clave} &
\textbf{Notas / respuestas} \\
\midrule

\textbf{¿Por qu\'e es obligatorio asistir a las audiencias?} &
Por dos motivos: uno pedag\'ogico (ver c\'omo se hace una audiencia real, con personas reales) y otro pr\'actico (los alumnos tienen informaci\'on del caso que el equipo docente no tiene, como si la parte pag\'o o si va a asistir). Es adem\'as un elemento para aprobar la materia. \\

\textbf{¿Qu\'e ocurre si no puede asistir todo el grupo?} &
Al menos una persona del grupo debe asistir. Se eval\'uan las situaciones laborales, personales o familiares, pero la ausencia total y no informada implica mala nota para todo el grupo; si nadie puede ir, debe avisarse al docente. \\

\textbf{¿Qu\'e actitud se observ\'o habitualmente en los profesionales que asisten a audiencias?} &
Apuro, desatenci\'on y exceso de audiencias; se dialoga en el pasillo sobre cu\'anto puede pagarse y se cierra la audiencia casi sin escuchar. La gente cree que es un tr\'amite. \\

\textbf{¿Qu\'e se espera del practicante respecto de su cliente durante la audiencia?} &
Estar cerca para poder moderar discretamente relatos extensos y evitar una mala impresi\'on ante el secretario. \\

\textbf{¿Qu\'e incentivo propone el docente?} &
Una tabla de posiciones de asistencia a tribunales, con foto como comprobante subida al grupo, y un premio. \\

\textbf{¿Qu\'e perfil profesional se busca formar?} &
El de un abogado pr\'actico, que tiene m\'as posibilidades de resolver un caso que uno ``voltero'' (apego a la teor\'ia). \\

\midrule

\multicolumn{2}{>{\raggedright\arraybackslash}X}{
\textbf{Resumen:}
La asistencia a audiencias es obligatoria por razones pedag\'ogicas y pr\'acticas, y cuenta para la aprobaci\'on del grupo. El docente destaca la actitud cercana y atenta hacia el cliente y el valor de la resoluci\'on pr\'actica de los conflictos por sobre el apego te\'orico.
} \\

\bottomrule
\end{tabularx}
\end{table}

% ============================================================
% CAPÍTULO 5
% ============================================================
\chapter[Compromiso social de la cl\'inica]{Cierre: compromiso social de la cl\'inica}

El docente introduce el cierre de la sesi\'on haciendo referencia al contexto institucional adverso que atraviesa la actividad acad\'emica, con medidas de fuerza (``paro'') en la facultad, y explica por qu\'e, pese a ello, el consultorio no adhiere a esas medidas: tiene una \textbf{funci\'on social} que considera esencial. All\'i se trabaja con personas que, si no son atendidas, gastaron lo que ten\'ian por el motivo que fuera y se van sin resolver su problema. Indica que es el \'unico motivo por el que no se adhiere y agrega que, en derecho, se vive en una ``burbuja'': se cree que no se forma parte de la universidad, pero se forma parte.

A modo de bienvenida y cierre, comparte con el grupo un fragmento de di\'alogo cinematogr\'afico entre un padre, radicado desde hace veinte a\~nos en el exterior, y su hijo, que discuten sobre la identidad nacional, el desarraigo y la posibilidad de cambio en Argentina. El docente aclara expresamente que la recomendaci\'on no constituye una toma de posici\'on partidaria, sino una invitaci\'on a reflexionar sobre el sentido del trabajo del consultorio. Los jueves se hacen recomendaciones; la pel\'icula es muy vieja y est\'a disponible ampliamente. Cada alumno decidir\'a de qu\'e lado est\'a y para qu\'e lado quiere trabajar.

El docente se\~nala que la pel\'icula tiene veintiocho a\~nos y que se sigue hablando de lo mismo. Expresa que se trabaja para la Argentina que quiere el hijo del di\'alogo, aunque est\'e en un papel sumiso o resignado, con esa esperanza, y que anhela que esto cambie en alg\'un momento. Aclara que no habla de pol\'itica, sino de lo que se vive todos los d\'ias en el consultorio.
% TODO: contenido incompleto o ambiguo en las notas originales (no se consigna el titulo de la pelicula; el guion del fragmento no se reproduce)

La sesi\'on finaliza con indicaciones organizativas menores: se solicita a los alumnos nuevos permanecer unos minutos m\'as para completar un registro de datos.

\begin{note}
El fragmento fílmico corresponde a una obra audiovisual protegida por derechos de autor; por ello no se reproduce su guion. Se deja constancia solo de su contenido tem\'atico (identidad, migraci\'on y esperanza de cambio en el pa\'is) para contextualizar la reflexi\'on del docente.
\end{note}

\section*{Cuadro Cornell: Compromiso social de la cl\'inica}
\addcontentsline{toc}{section}{Cuadro Cornell: Compromiso social de la cl\'inica}

\begin{table}[H]
\centering
\small
\begin{tabularx}{\textwidth}{
    >{\raggedright\arraybackslash}p{0.28\textwidth}
    >{\raggedright\arraybackslash}X
}
\toprule
\textbf{Preguntas / palabras clave} &
\textbf{Notas / respuestas} \\
\midrule

\textbf{¿Por qu\'e el consultorio no adhiere a las medidas de fuerza?} &
Porque tiene una funci\'on social esencial: las personas consultantes, si no son atendidas, se van sin resolver su problema tras haber gastado lo que ten\'ian. \\

\textbf{¿Qu\'e quiere decir que en derecho se vive en una ``burbuja''?} &
Que se cree no formar parte de la universidad, cuando s\'i se forma parte. \\

\textbf{¿Qu\'e fragmento comparte el docente y con qu\'e finalidad?} &
Un di\'alogo cinematogr\'afico entre un padre radicado veinte a\~nos en el exterior y su hijo sobre identidad nacional, desarraigo y posibilidad de cambio; no es una toma de posici\'on partidaria sino una invitaci\'on a reflexionar sobre el sentido del trabajo del consultorio. \\

\textbf{¿Qu\'e esperanza expresa el docente?} &
Que algo cambie en alg\'un momento; aclara que habla de lo que se vive cada d\'ia en el consultorio y no de pol\'itica. \\

\midrule

\multicolumn{2}{>{\raggedright\arraybackslash}X}{
\textbf{Resumen:}
El consultorio mantiene su actividad pese a las medidas de fuerza institucionales porque su funci\'on social es esencial para quienes no pueden esperar. El cierre invita a cada alumno a reflexionar sobre para qui\'en y para qu\'e quiere trabajar.
} \\

\bottomrule
\end{tabularx}
\end{table}

\end{document}
